\begin{textsample}{POS Dim 3 – df – Score 15.00 – df/df\_inf\_3.txt}  \label{ex:f3_pos_126}
3.
O DISTRITO FEDERAL E SUAS \textbf{CRIANÇAS} : um olhar à diversidade cultural das \textbf{infâncias} Dando continuidade à discussão da territorialidade, indaga -se : que \textbf{crianças} e \textbf{infâncias} estão presentes no Distrito Federal?
Como tem se visto, os conceitos que identificam a \textbf{infância} se constituíram ao longo da história até se \textbf{depararem} com a \textbf{criança} definida como sujeito histórico de direitos, atuante e protagonista na constituição de sua identidade pessoal e coletiva.
Mediante suas interações, relações e práticas cotidianas, a \textbf{criança} utiliza o \textbf{brincar}, a imaginação, a fantasia, a observação, as narrativas, os questionamentos, “ experimenta, aprende e constrói sentidos sobre a natureza e a sociedade, produzindo cultura ” ( BRASIL, 2010a, p. 12 ).
Partindo dessa compreensão de \textbf{criança}, cabe observar suas \textbf{infâncias} e seus percursos como produtoras de cultura, pois há inegável diversidade cultural brasileira que se reflete na composição do Distrito Federal, dadas suas peculiaridades que comportam tanto os modos de viver das \textbf{crianças} do campo, indígenas, quilombolas e migrantes do território nacional.
O trabalho educativo nas \textbf{instituições} que ofertam Educação \textbf{Infantil} pressupõe a ampliação do olhar voltado às \textbf{infâncias} constituídas historicamente no território distrital, pois : > \textbf{Crianças} e \textbf{infâncias} são marcadas por conceitos constituídos social e culturalmente.
O modo como são percebidas e compreendidas interfere, direta e indiretamente, na organização do trabalho pedagógico a ser realizado nas \textbf{instituições} educativas para a primeira \textbf{infância} ( DISTRITO FEDERAL, 2018a, p. 21 ).
Tal representatividade cultural presente no Distrito Federal abarca, também, os contextos das \textbf{crianças} estrangeiras, reiterando assim, que “ dentro das \textbf{instituições} que ofertam Educação \textbf{Infantil} na SEEDF, pressupõe que não podemos nos restringir a pensar uma única \textbf{criança}, uma única \textbf{infância} e, sim, às inúmeras \textbf{infâncias} e \textbf{crianças} que se fazem presentes no contexto educativo e coletivo ” ( DISTRITO FEDERAL, 2018a, p. 21 ).
Nesse sentido, observa -se que : > A comunidade escolar está cada vez mais diversificada, hoje temos em nossa rede de ensino, \textbf{crianças} indígenas, quilombolas, do campo, entre outras, envolvidas em um mar de tecnologias, que podem ter ou não mais ou menos influência no seu cotidiano.
Todas essas diferentes \textbf{crianças}, com especificidades distintas, precisam ser consideradas na prática educativa ( DISTRITO FEDERAL, 2018a, p. 21 ).
A exemplo do eixo \textbf{brincadeira}, evidenciado pelo Guia da VI Plenarinha ( DISTRITO FEDERAL, 2018a ), torna -se importante destacar esta pluralidade \textbf{infantil} no trabalho educativo, nos diversos campos de experiências elencados neste Currículo.
A prática docente, permeando os campos de experiências, os direitos e objetivos de aprendizagem e desenvolvimento da \textbf{criança}, deve compreender e considerar, em sua intencionalidade educativa, as variáveis que constituem as \textbf{infâncias} presentes no Distrito Federal alinhadas às práticas socioculturais da atualidade.
Mesmo dentro de um grupo étnico existem muitas diferenças, como entre os indígenas.
As \textbf{crianças} indígenas têm sua \textbf{infância} diferenciada em cada etnia, formando em sua aldeia seu modo de vida social particular.
Portanto, deve -se considerar, na \textbf{instituição} que oferta Educação \textbf{Infantil}, os momentos de trocas de experiências, em que as \textbf{crianças} compartilhem sua cultura, aprendendo umas com as outras, enriquecendo seu repertório cultural e de seus \textbf{pares}.
É importante mencionar as experiências que alguns grupos sociais possuem em meio à coletividade, com os anciãos, com o território, com os recursos naturais, as tradições, os costumes, entre outras, que podem ser compartilhadas com os que não usufruem desse contexto, inclusive em atividades realizadas nos espaços fora dos muros da \textbf{instituição} educativa. > Lembrando que as tradições não são isoladas, elas se movimentam com o fluxo do processo histórico, possibilitando o resgate e a continuidade de tradições, dialogando com a sociedade atual, em uma dinâmica de resistência e transformação ( DISTRITO FEDERAL, 2018a, p. 23 ).
Destaca -se que a intencionalidade do trabalho educativo com \textbf{crianças} das mais diversas culturas deve estabelecer vínculos com seus valores culturais, sociais, históricos e econômicos de suas comunidades, onde a \textbf{instituição} que oferta Educação \textbf{Infantil} se estabelece como “ um espaço de diálogo entre o conhecimento escolar e a realidade social das \textbf{crianças}, valorizando o desenvolvimento sustentável, o trabalho, a cultura, a luta pelo direito à terra e ao território ” ( DISTRITO FEDERAL, 2018a, p. 23 ).
No cotidiano da Educação \textbf{Infantil}, o docente deverá propiciar momentos de \textbf{escuta} e \textbf{rodas} de conversa com vistas a identificar as características culturais individuais das \textbf{crianças}.
Em geral, as \textbf{crianças} quilombolas, indígenas, bem como as do campo, convivem nos espaços naturais disponíveis e se desenvolvem nas atividades realizadas nesse contexto. > As escolas que recebem \textbf{crianças} indígenas e quilombolas ou de outros grupos sociais, devem olhar as especificidades dessas \textbf{crianças}, respeitando suas histórias de vida, suas comunidades, suas tradições e \textbf{convidando} -as a compartilhar suas experiências, possibilitando espaços e tempos de troca, para que assim, todas as \textbf{crianças} possam ampliar suas experiências, saberes e conhecimentos ( DISTRITO FEDERAL, 2018a, p. 24 ).
Em relação às \textbf{crianças} do campo, vale destacar que as \textbf{rotinas} ambientais, \textbf{sensoriais}, afetivas, culturais e sociais envolvem particularidades voltadas aos ciclos de produção, à sazonalidade de plantio e de colheita, de estiagem das águas, das épocas de reprodução dos peixes, das aves e de outros animais, de idas e vindas de aves e de outros bichos ( SILVA ; PASUCH, 2010 ).
Outro aspecto importante a considerar na Educação \textbf{Infantil} é o desenvolvimento de uma educação que promova a igualdade racial, no sentido de apresentar às \textbf{crianças} a realidade existente e provocar reflexões sobre a diversidade humana e o respeito a essa diversidade, isso em relação a todas as raças e etnias que constituem a humanidade.
Se uma \textbf{criança} negra se sente bem com o seu corpo, seu rosto e seus cabelos, e uma \textbf{criança} branca também se sente bem consigo mesma, pode haver respeito e aceitação entre elas.
Essa é a importância do trabalho com a promoção da igualdade racial nesta etapa.
Se houver uma intervenção qualificada e que não ignore a “ raça ” como um componente importante no processo de construção da identidade da \textbf{criança}, teremos outra história sendo construída ( BRASIL, 2012, p. 9 ).
Como referendado acima, o respeito é essencial no processo educativo e, nesse sentido, destaca -se as \textbf{crianças} indígenas, estrangeiras e/ou refugiadas que precisam de atenção em relação à sua inserção em uma nova cultura que possui muitas diferenças, dentre elas, a língua.
Em muitas etnias indígenas, as \textbf{crianças} dialogam usando apenas sua língua materna, situação que também ocorre com as \textbf{crianças} de outros países.
As \textbf{instituições} que atendem à Educação \textbf{Infantil} precisam se organizar para \textbf{acolher} essas \textbf{crianças}, de forma que tenham um olhar e uma \textbf{escuta} sensível às suas necessidades, buscando estratégias de comunicação e de inserção no coletivo das próprias \textbf{crianças}, por meio de inúmeras formas de expressão que podem ser vivenciadas pelos humanos.
Cabe ainda destacar o esforço no sentido de estabelecer comunicação entre a \textbf{instituição} educativa, a \textbf{criança} e sua família e/ou responsáveis, focando nos aspectos afetivos e cognitivos, bem como motores, \textbf{sensoriais} e sociais, imbricados nas relações educativas.
Portanto, faz -se necessário desenvolver um olhar e uma \textbf{escuta} \textbf{atenta} à cultura, respeitando histórias e modos de vida e de estar no mundo da \textbf{criança}, bem como sua formação identitária nas relações que estabelece com sua cultura.
Dessa maneira, a \textbf{instituição} que oferta Educação \textbf{Infantil} deve proporcionar ocasiões de trocas de vivências e experiências entre as diversas \textbf{infâncias} existentes em seus espaços educativos, ampliando as possibilidades de desenvolvimento de cada \textbf{criança} como sujeito que se constitui também nesse espaço social.

% matched lemmas: acolher, atento, brincadeira, brincar, convidar, criança, deparar, escuta, infantil, infância, instituição, par, roda, rotina, sensorial
\end{textsample}
